\section{Exact common predictions}\label{sec:preparation}

Initial peeling gives closed pins and spreading leaves. Each marked
effective profile has total component rank at most the initial pin budget
\(e\). We will use this bound directly.

\subsection{An exact prediction alternative}

A common predictor at endpoint \(i\) has the form
\begin{equation}\label{eq:linear-predictor}
 f_i(l,q,k)=\varphi_{i,l}(q,L_{i,l}k).
\end{equation}
For each retained unit there are effective marks \(\lambda_i\), bits
\(\alpha_i\), and bounded sets of excluded labels. The prediction is
\begin{equation}\label{eq:common-prediction}
 p_i(A)=T(\lambda_i,A)+\alpha_i a(A)+f_i(l,q,k),
\end{equation}
exactly, at every generic input for each allowed label. The finite alternative
asks whether such predictors, with the required linear coordinate bound,
exist on unit mass at least \(10^{-6}\). The alternative without a
prediction is treated in \Cref{sec:status-span}. There, relations on
possible statuses will produce these identities directly.

An exact identity on generic inputs also holds at shared points: set the
ordinary coordinates to zero. This preserves the atom, its key and its
tester value. Thus no separate prediction assumption on shared inputs is
needed.

\subsection{Free coordinates and exact sections}

Pull all common coordinates at one endpoint back to its array of component
keys and take their span \(H\). This identifies duplicated coordinates.
The raw frame count bounds the probability that a fixed nonzero row
vanishes on the free \(Z\)-coordinates by \(8\cdot2^{-n}\).
After a union over rows and labels, a prepared marginal bounded by
\(2^{n/32}\mu\) satisfies
\begin{equation}\label{eq:prepared-rank}
 \PP(\text{some label fails full rank on }H)
 \le8\cdot2^{\dim H+b-31n/32}.
\end{equation}
The row bound gives \(\dim H\le n/4\). The exceptional mixed mass is
therefore exponentially small. It is discarded before restricting leaf
pairs or imposing compatibility.

Fix a unit outside this exception. At a fixed label the common coordinates
are \(c=AZ+By+c_0\), where \(A\) is onto and \(y\) contains the other
variables. A right inverse \(j\) gives
\[
 Z=j(c-By-c_0)+z',\qquad z'\in\ker A.
\]
Now \(c,y,z'\) are independent and uniform. The tester \(q=q(y)\)
is independent of \(c\), and both its values have probability at least
\(1/4\).

Fix one retained unit and one allowed label. After moving \(f_i\)
to the other side, the prediction reads
\[
 R_l(c,y,z')=\varphi_l(q(y),c),
\]
where \(R_l\) is quadratic. Fix any \((y,z')\) with \(q=0\).
Restriction to this section gives a quadratic polynomial
\(Q_{l,0}(c)=\varphi_l(0,c)\). A section with \(q=1\) gives
\(Q_{l,1}(c)=\varphi_l(1,c)\). These are exact equalities.

Choose one disjoint tester pair \(y_a,y_b\) in \(q\). Hold \(c,z'\)
fixed and take the second difference in those two directions. On the
left this is constant, since \(R_l\) is quadratic. On the right it is
\(Q_{l,0}(c)+Q_{l,1}(c)\). Thus the difference is a constant
\(\gamma_l\), and
\[
 \varphi_l(q,c)=Q_{l,0}(c)+\gamma_l q.
\]

The function \(\varphi_l\) is common to the population. Substituting
this equality therefore gives the same identity at every other retained
unit, at every allowed label. Fix a base point. Its difference is a
polynomial of degree at most 2 in the selector bits. A nonzero Boolean
polynomial of degree at most \(d\) is nonzero on at least \(2^{-d}\)
of its inputs, by induction on the variables. Fewer than a quarter of
the labels are excluded, so the identity extends to every label. It
holds at every base point, including shared points.

We have obtained common polynomials \(Q_{i,l}\) and constant bits
\(\gamma_{i,l}\) such that
\begin{equation}\label{eq:fitted-prediction}
 p_i(A)=T(\lambda_i,A)+\alpha_i a(A)
       +Q_{i,l}(L_{i,l}K_i(A))+\gamma_{i,l}q(A).
\end{equation}
On these sections the tester terms are constant. The other terms
have a matrix representative of rank at most \(r_{\rm scr}\).
Matrix factorization shows that each \(Q_{i,l}\) uses at most
\(R_*=2r_{\rm scr}\) linear coordinates. This bound is independent of
the original partition size.

\subsection{The tester slope is zero}

Put
\[
 v_i(t)=a(\lambda_i)+b_t(\lambda_i)+\alpha_i.
\]
Fix a common cell in \eqref{eq:fitted-prediction}. Solving for \(Z\)
leaves every tester variable free. The ordinary tester part is
\(\sum_t(v_i(t)+\gamma_{i,l})\eta_{l,t}\); the remaining part has
matrix rank at most \(r_{\rm scr}\). Comparing their second differences
gives the conservative bound used in Lean,
\begin{equation}\label{eq:tester-support}
 r_0|\supp(v_i+\gamma_{i,l}\mathbf1)|\le2r_{\rm scr}.
\end{equation}
One can read this comparison on disjoint tester pairs: their second
differences contain an identity matrix, whereas the polar matrix of the
remaining expression has rank at most \(2r_{\rm scr}\).

If two tags had different slopes, the vectors in
\eqref{eq:tester-support} would be complements. The inequality
\(4r_{\rm scr}<gr_0\) rules this out. Thus the slope is a single bit
\(\gamma_i\). Use the same shared point at every tag. Its atoms sum to
zero, so their prediction values sum to zero. Their joint common cell
is independent of \(q\). Since \(g\) is odd, changing \(q\) changes
that sum by \(\gamma_i\), which must therefore be zero.

Absorb the linear coordinate maps into the notation for \(Q\). We have
proved the exact identity
\begin{equation}\label{eq:exact-quadratic-prediction}
 p_i(A)=T(\lambda_i,A)+\alpha_i a(A)+Q_{i,l}(K_i(A))
\end{equation}
for every label and every base point at every retained unit.

\subsection{Finite classes and consistency}

The prediction classes are literal equalities. An endpoint is inactive
when \(\alpha_i=1\) and every \(Q_{i,l}=0\). Two active endpoints have
the same class when their bits and polynomials agree. The other possibilities
are distinct classes.

The two predictor values depend on at most \(2R_*\) linear coordinates.
Every value pair attained at a key contains a full fiber of these
coordinates. That fiber has uniform mass at least \(2^{-2R_*}\).
The reference Fourier bound gives it mass at least \(2^{-2R_*}/2\)
on each diagonal slice for all sufficiently large \(n\). Thus both
slices contain every attained pair. The scalar argument can test them all.

For an inactive endpoint, \eqref{eq:exact-quadratic-prediction} gives
\begin{equation}\label{eq:exact-zero-class}
 p_i=T(\lambda_i,\cdot)+a\quad\text{on }\cX,
\end{equation}
because atoms span \(\cX\).

Suppose now that \(\Delta=0\), so \(U_1\lambda_1=U_2\lambda_2\).
Sum the equality of component tester values
\eqref{eq:common-image-tester} along \(M_t\). Privacy of the ordinary
blocks gives
\[
 v_1+v_2=(\alpha_1+\alpha_2)\mathbf1+o(\lambda_1)+o(\lambda_2).
\]
The marks have total rank at most \(e\), so each ordinary support has
size at most \(e\) by \Cref{lem:ordinary-support}. If the bits differed,
the all-one vector would be a sum of four vectors whose total support is
at most \(4r_{\rm scr}/r_0+2e<g\). This is impossible. Thus
\(\alpha_1=\alpha_2\).

\subsection{The prepared alternatives}

Fix \(c=a(\lambda_1)+a(\lambda_2)\). Write
\(v_i=u_{B_i}(\Delta_A)\), \(w_i=u_{A_i}(\Delta_B)\) in the following
table; these are new scalar symbols, not the tester vectors above.
\begin{center}
\begin{tabular}{ll}
\toprule
Prediction class & Required relation\\
\midrule
Both inactive & \(v_1=v_2=w_1=w_2=c\)\\
Only endpoint \(i\) active & \(v_i+w_i=1+\alpha_i\)\\
Equal active pair & \(v_1+v_2+w_1+w_2=0\)\\
No further relation & No condition\\
\bottomrule
\end{tabular}
\end{center}
The dual calculation in \Cref{sec:status-span} proves this table.
When both endpoints are inactive, \eqref{eq:exact-zero-class} gives
\((u_{A_1}+u_{A_2})(\Delta_A)=0\). If also \(\Delta_A=0\), then
\(c=0\); otherwise the two marks would witness an internal hole.

Distinct active predictions permit every phase. When \(\Delta=0\),
consistency gives compatibility. For the two reciprocal zero parities,
\eqref{eq:rank-compatibility} gives compatible mass of order \(N^{-2}\).
If ordinary trimming leaves a conditional probability of at least
\(1/2\) wherever compatibility occurs, and its unsafe share is at most
\(0.001\), this settles the case. When both predictions are inactive and the target
is zero, the four-zero count above applies. In each of these rank
arguments, if the retained offsets span a space within the label budget of
\Cref{lem:injection}, refine by the labels just described. Each
compatible leaf pair now has conditional probability one, and
\Cref{lem:injection} controls its unsafe share in the aggregate.
For the two prescribed-one cases, the same refinement works if every
set of mass greater than \(1/2\) contains a compatible pair. Here is one
way to check that condition when the offsets are nonzero. Suppose the
offsets on each set without a compatible pair span at most \(N/2\)
dimensions. A fixed nonzero offset has a nonzero component whose column
space meets one of the two primal images. Indeed, quotient the first
factor by the first image and the second factor by the second image.
The offset maps to zero. If both quotients were injective on its column
space, their tensor product would be injective there as well, since
the component is symmetric. The frame count bounds the probability of
this intersection. Summing over the at most \(2^{N/2}\) offsets of the
given subspace gives mass at most \(2^{-N/8}\) for all sufficiently
large \(n\). Such a set cannot have mass greater than \(1/2\).
For the remaining laws use \Cref{thm:phase-alternative}. Its mixed case
gives positive compatible mass. Its small cover gives the required
label budget and the singleton estimates for prescribed bits.

The mass bound is elementary. Suppose a first restriction has mass at
least \(a>0\), and the next has conditional mass at least \(b>0\).
If trimming deletes at most \(\delta\le ab/2\) of the original mass,
the remainder has mass at least \(ab/2\). For any pair event, its
probability after normalizing decreases by at most \(2\delta/(ab)\).
Here \(a,b\) are fixed constants and every trimming error is exponentially
small. This gives the claimed bound through both preparation branches.

% lax begin Lax793645.PreparedUnits
\begin{proposition}\label{prop:status-preparation}
There exist \(a>0\) and \(n_0\) such that every admissible law at
every \(n\ge n_0\) admits a prepared population of mass at least \(a\).
It satisfies \eqref{eq:final-leaf-minentropy}. In the absence of common
predictions more than \(0.99\) of pairs are safe. Otherwise the
predictions are exact. Either safe compatible pairs have mass at least
\(a\), or compatible pairs have mass at least \(a/N^4\), their unsafe
share is at most \(0.001\), and each positive conditional compatibility
probability on a leaf pair is at least \(a\).
\end{proposition}
% lax end

% lax begin Lax793645Proofs.Manuscript.PreparedUnits.one_finite_preparation
\begin{proof}
For uniformity, write the possible lower bounds as \(a_k=1/(k+1)\).
Within each of the finitely many cases the preceding estimates hold
eventually on the dimensions where that case occurs. Take the largest
required index \(k\) across the cases. This works for any sequence of input
laws. To make \(k\) uniform in the law, consider the indices \(j\le n\)
for which some input fails at \(a_j\). Choose a failing input for the
largest such \(j\).
If no input fails there is nothing to prove at that dimension. The
bound \(a_k\) for these chosen inputs works eventually for every input:
a failure at index \(k\le n\) would give a chosen failure at an index
\(j\ge k\), whereas success at \(k\) implies success at \(j\).
\end{proof}
% lax end
